\documentclass{article}
\usepackage{iclr2027_conference}
\usepackage{fontspec}
\setmainfont{texgyretermes-regular.otf}[Path=fonts/,BoldFont=texgyretermes-bold.otf,ItalicFont=texgyretermes-italic.otf,BoldItalicFont=texgyretermes-bolditalic.otf]
\usepackage{amsmath,amssymb,booktabs,longtable,xurl}
\usepackage[colorlinks=true,allcolors=blue]{hyperref}
\usepackage{microtype}
\hypersetup{pdftitle={BioCoLoop: Statistical Supplement}}
\begin{document}
\title{BioCoLoop: Statistical Supplement}
\author{Anonymous}
\maketitle
\renewcommand{\thetable}{S.\arabic{table}}
\renewcommand{\theHtable}{S.\arabic{table}}

This supplement accompanies \emph{BioCoLoop: Collaborative Agentic Research for Biological Model Improvement}. It reports all per-seed comparisons from the two complete uncertainty tables, without selecting results by direction or significance. Tables S.1 and S.2 retain their original numerical records. Mean scores, the comparison of the one- and ten-laboratory conditions and the decomposition of effects remain in Appendix A of the manuscript.

Intervals are paired cluster-bootstrap 95\% intervals from 1,000 fixed-seed resamples, conditional on the selected trained predictors. Drug identity defines clusters for random and unseen-drug interaction tests, protein sequence for unseen-protein tests, and compound or intervention identity for proteomic or cell tests. Each contrast uses the same sampled clusters for both methods. Seeds 42--44 use fixed data partitions. These intervals measure test-sample uncertainty, separately from training-seed variation.

\begin{longtable}{llrrl}
\caption{Per-seed paired cluster-bootstrap 95\% intervals for the committed model versions, in percentage points. These retrospective intervals condition on fixed trained models and do not measure across-seed uncertainty. DTI intervals, when available, are a post-hoc supplement using frozen predictions; random-endpoint rows use drug clusters.}\label{tab:strong-uncertainty}\\
\toprule Task & Contrast & Seed & Difference & 95\% interval \\ \midrule\endfirsthead
\toprule Task & Contrast & Seed & Difference & 95\% interval \\ \midrule\endhead
DTI & Loop $-$ fixed (10 labs) & 42 & +0.33 & [-0.06, +0.69] \\
DTI & Loop $-$ direct (1 lab) & 42 & +0.00 & [+0.00, +0.00] \\
DTI & Loop $-$ direct (10 labs) & 42 & +0.00 & [+0.00, +0.00] \\
DTI & 10 $-$ 1 labs (both fixed) & 42 & +12.51 & [+11.35, +13.71] \\
DTI & 10 $-$ 1 labs (both direct) & 42 & +5.78 & [+5.02, +6.57] \\
DTI & 10 $-$ 1 labs (both loop) & 42 & +5.78 & [+5.02, +6.57] \\
\midrule
DTI & Loop $-$ fixed (10 labs) & 43 & +0.54 & [-0.08, +1.11] \\
DTI & Loop $-$ direct (1 lab) & 43 & +0.00 & [+0.00, +0.00] \\
DTI & Loop $-$ direct (10 labs) & 43 & +0.00 & [+0.00, +0.00] \\
DTI & 10 $-$ 1 labs (both fixed) & 43 & +7.80 & [+6.82, +8.82] \\
DTI & 10 $-$ 1 labs (both direct) & 43 & +1.65 & [+0.90, +2.38] \\
DTI & 10 $-$ 1 labs (both loop) & 43 & +1.65 & [+0.90, +2.38] \\
\midrule
DTI & Loop $-$ fixed (10 labs) & 44 & +5.01 & [+4.25, +5.71] \\
DTI & Loop $-$ direct (1 lab) & 44 & +0.00 & [+0.00, +0.00] \\
DTI & Loop $-$ direct (10 labs) & 44 & +0.00 & [+0.00, +0.00] \\
DTI & 10 $-$ 1 labs (both fixed) & 44 & +6.60 & [+5.61, +7.74] \\
DTI & 10 $-$ 1 labs (both direct) & 44 & +4.17 & [+3.45, +4.93] \\
DTI & 10 $-$ 1 labs (both loop) & 44 & +4.17 & [+3.45, +4.93] \\
\midrule
Proteomics & Loop $-$ fixed (10 labs) & 42 & +0.00 & [+0.00, +0.00] \\
Proteomics & Loop $-$ direct (1 lab) & 42 & +0.00 & [+0.00, +0.00] \\
Proteomics & Loop $-$ direct (10 labs) & 42 & +0.00 & [+0.00, +0.00] \\
Proteomics & 10 $-$ 1 labs (both fixed) & 42 & +13.16 & [-0.39, +26.02] \\
Proteomics & 10 $-$ 1 labs (both direct) & 42 & +16.91 & [+3.37, +31.02] \\
Proteomics & 10 $-$ 1 labs (both loop) & 42 & +16.91 & [+3.37, +31.02] \\
\midrule
Proteomics & Loop $-$ fixed (10 labs) & 43 & +0.00 & [+0.00, +0.00] \\
Proteomics & Loop $-$ direct (1 lab) & 43 & +0.00 & [+0.00, +0.00] \\
Proteomics & Loop $-$ direct (10 labs) & 43 & +0.00 & [+0.00, +0.00] \\
Proteomics & 10 $-$ 1 labs (both fixed) & 43 & +15.98 & [+1.29, +29.00] \\
Proteomics & 10 $-$ 1 labs (both direct) & 43 & +17.03 & [+3.34, +31.44] \\
Proteomics & 10 $-$ 1 labs (both loop) & 43 & +17.03 & [+3.34, +31.44] \\
\midrule
Proteomics & Loop $-$ fixed (10 labs) & 44 & -0.24 & [-3.40, +2.33] \\
Proteomics & Loop $-$ direct (1 lab) & 44 & +0.00 & [+0.00, +0.00] \\
Proteomics & Loop $-$ direct (10 labs) & 44 & +0.00 & [+0.00, +0.00] \\
Proteomics & 10 $-$ 1 labs (both fixed) & 44 & +6.93 & [-3.43, +15.60] \\
Proteomics & 10 $-$ 1 labs (both direct) & 44 & +16.63 & [+3.02, +31.08] \\
Proteomics & 10 $-$ 1 labs (both loop) & 44 & +16.63 & [+3.02, +31.08] \\
\midrule
VCC & Loop $-$ fixed (10 labs) & 42 & +0.00 & [+0.00, +0.00] \\
VCC & Loop $-$ direct (1 lab) & 42 & +0.00 & [+0.00, +0.00] \\
VCC & Loop $-$ direct (10 labs) & 42 & +0.00 & [+0.00, +0.00] \\
VCC & 10 $-$ 1 labs (both fixed) & 42 & +1.70 & [-5.36, +8.58] \\
VCC & 10 $-$ 1 labs (both direct) & 42 & +9.90 & [-2.88, +23.04] \\
VCC & 10 $-$ 1 labs (both loop) & 42 & +9.90 & [-2.88, +23.04] \\
\midrule
VCC & Loop $-$ fixed (10 labs) & 43 & +0.00 & [+0.00, +0.00] \\
VCC & Loop $-$ direct (1 lab) & 43 & +4.83 & [-5.50, +16.49] \\
VCC & Loop $-$ direct (10 labs) & 43 & +0.00 & [+0.00, +0.00] \\
VCC & 10 $-$ 1 labs (both fixed) & 43 & +1.56 & [-5.76, +8.58] \\
VCC & 10 $-$ 1 labs (both direct) & 43 & +15.08 & [+2.94, +27.78] \\
VCC & 10 $-$ 1 labs (both loop) & 43 & +10.24 & [-2.74, +23.49] \\
\midrule
VCC & Loop $-$ fixed (10 labs) & 44 & +0.00 & [+0.00, +0.00] \\
VCC & Loop $-$ direct (1 lab) & 44 & -7.44 & [-17.23, +1.33] \\
VCC & Loop $-$ direct (10 labs) & 44 & +0.00 & [+0.00, +0.00] \\
VCC & 10 $-$ 1 labs (both fixed) & 44 & +1.56 & [-5.76, +8.58] \\
VCC & 10 $-$ 1 labs (both direct) & 44 & +6.38 & [-4.38, +19.13] \\
VCC & 10 $-$ 1 labs (both loop) & 44 & +13.82 & [+2.54, +26.34] \\
\midrule
Norman & Loop $-$ fixed (10 labs) & 42 & +7.26 & [-9.41, +25.41] \\
Norman & Loop $-$ direct (1 lab) & 42 & +0.00 & [+0.00, +0.00] \\
Norman & Loop $-$ direct (10 labs) & 42 & +0.00 & [+0.00, +0.00] \\
Norman & 10 $-$ 1 labs (both fixed) & 42 & -10.07 & [-21.19, +0.00] \\
Norman & 10 $-$ 1 labs (both direct) & 42 & +8.59 & [-8.30, +26.00] \\
Norman & 10 $-$ 1 labs (both loop) & 42 & +8.59 & [-8.30, +26.00] \\
\midrule
Norman & Loop $-$ fixed (10 labs) & 43 & +6.37 & [-11.12, +24.81] \\
Norman & Loop $-$ direct (1 lab) & 43 & +0.00 & [+0.00, +0.00] \\
Norman & Loop $-$ direct (10 labs) & 43 & +0.00 & [+0.00, +0.00] \\
Norman & 10 $-$ 1 labs (both fixed) & 43 & -10.07 & [-21.19, +0.00] \\
Norman & 10 $-$ 1 labs (both direct) & 43 & +7.70 & [-9.33, +25.85] \\
Norman & 10 $-$ 1 labs (both loop) & 43 & +7.70 & [-9.33, +25.85] \\
\midrule
Norman & Loop $-$ fixed (10 labs) & 44 & +6.67 & [-10.37, +24.81] \\
Norman & Loop $-$ direct (1 lab) & 44 & +0.00 & [+0.00, +0.00] \\
Norman & Loop $-$ direct (10 labs) & 44 & +0.00 & [+0.00, +0.00] \\
Norman & 10 $-$ 1 labs (both fixed) & 44 & -10.07 & [-21.19, +0.00] \\
Norman & 10 $-$ 1 labs (both direct) & 44 & +8.00 & [-8.52, +25.93] \\
Norman & 10 $-$ 1 labs (both loop) & 44 & +8.00 & [-8.52, +25.93] \\
\midrule
Tahoe & Loop $-$ fixed (10 labs) & 42 & +0.00 & [+0.00, +0.00] \\
Tahoe & Loop $-$ direct (1 lab) & 42 & +0.00 & [+0.00, +0.00] \\
Tahoe & Loop $-$ direct (10 labs) & 42 & +0.00 & [+0.00, +0.00] \\
Tahoe & 10 $-$ 1 labs (both fixed) & 42 & +4.48 & [-5.97, +14.93] \\
Tahoe & 10 $-$ 1 labs (both direct) & 42 & +10.45 & [-1.49, +22.39] \\
Tahoe & 10 $-$ 1 labs (both loop) & 42 & +10.45 & [-1.49, +22.39] \\
\midrule
Tahoe & Loop $-$ fixed (10 labs) & 43 & -2.99 & [-7.46, +0.00] \\
Tahoe & Loop $-$ direct (1 lab) & 43 & +0.00 & [+0.00, +0.00] \\
Tahoe & Loop $-$ direct (10 labs) & 43 & +0.00 & [+0.00, +0.00] \\
Tahoe & 10 $-$ 1 labs (both fixed) & 43 & +7.46 & [-4.48, +17.95] \\
Tahoe & 10 $-$ 1 labs (both direct) & 43 & +13.43 & [+1.49, +23.88] \\
Tahoe & 10 $-$ 1 labs (both loop) & 43 & +13.43 & [+1.49, +23.88] \\
\midrule
Tahoe & Loop $-$ fixed (10 labs) & 44 & -7.46 & [-19.40, +4.48] \\
Tahoe & Loop $-$ direct (1 lab) & 44 & +0.00 & [+0.00, +0.00] \\
Tahoe & Loop $-$ direct (10 labs) & 44 & +0.00 & [+0.00, +0.00] \\
Tahoe & 10 $-$ 1 labs (both fixed) & 44 & +7.46 & [-4.48, +17.95] \\
Tahoe & 10 $-$ 1 labs (both direct) & 44 & -2.99 & [-17.91, +11.94] \\
Tahoe & 10 $-$ 1 labs (both loop) & 44 & -2.99 & [-17.91, +11.94] \\
\midrule
\bottomrule
\end{longtable}

\clearpage
\begingroup\scriptsize\setlength{\tabcolsep}{3pt}
\begin{longtable}{llrrrr}
\caption{Post-hoc TAPB uncertainty supplement from frozen prediction files. All three endpoints, six arms and fixed contrasts are retained. Entries are paired AUROC/AP differences and percentile 95\% intervals in percentage points from 1,000 requested cluster draws (RNG seed 20260918). Random and unseen-drug endpoints resample drug clusters; unseen-protein resamples protein clusters. Intervals condition on fixed trained models and one clustering axis; they are not two-way, across-seed, multiplicity-adjusted, or prospective confirmatory intervals.}\label{tab:strong-dti-uncertainty}\\
\toprule Endpoint / seed & Contrast & AUROC $\Delta$ & 95\% interval & AP $\Delta$ & 95\% interval \\ \midrule\endfirsthead
\toprule Endpoint / seed & Contrast & AUROC $\Delta$ & 95\% interval & AP $\Delta$ & 95\% interval \\ \midrule\endhead
Random / 42 & Loop $-$ fixed (10 labs) & +0.33 & [-0.06, +0.69] & +1.40 & [+0.76, +2.05] \\
Random / 42 & Loop $-$ direct (1 lab) & +0.00 & [+0.00, +0.00] & +0.00 & [+0.00, +0.00] \\
Random / 42 & Loop $-$ direct (10 labs) & +0.00 & [+0.00, +0.00] & +0.00 & [+0.00, +0.00] \\
Random / 42 & 10 $-$ 1 labs (fixed) & +12.51 & [+11.35, +13.71] & +10.56 & [+9.10, +12.11] \\
Random / 42 & 10 $-$ 1 labs (direct) & +5.78 & [+5.02, +6.57] & +8.29 & [+7.01, +9.56] \\
Random / 42 & 10 $-$ 1 labs (loop) & +5.78 & [+5.02, +6.57] & +8.29 & [+7.01, +9.56] \\
Random / 42 & Our harness $-$ fixed (1 lab) & +12.84 & [+11.82, +14.01] & +11.96 & [+10.56, +13.45] \\
Random / 42 & Our harness $-$ direct (1 lab) & +5.78 & [+5.02, +6.57] & +8.29 & [+7.01, +9.56] \\
\midrule
Unseen drug / 42 & Loop $-$ fixed (10 labs) & +0.10 & [-0.47, +0.67] & +1.06 & [+0.42, +1.66] \\
Unseen drug / 42 & Loop $-$ direct (1 lab) & +0.00 & [+0.00, +0.00] & +0.00 & [+0.00, +0.00] \\
Unseen drug / 42 & Loop $-$ direct (10 labs) & +0.00 & [+0.00, +0.00] & +0.00 & [+0.00, +0.00] \\
Unseen drug / 42 & 10 $-$ 1 labs (fixed) & +12.30 & [+10.77, +13.98] & +7.17 & [+5.82, +8.72] \\
Unseen drug / 42 & 10 $-$ 1 labs (direct) & +6.38 & [+5.34, +7.53] & +7.55 & [+6.30, +8.92] \\
Unseen drug / 42 & 10 $-$ 1 labs (loop) & +6.38 & [+5.34, +7.53] & +7.55 & [+6.30, +8.92] \\
Unseen drug / 42 & Our harness $-$ fixed (1 lab) & +12.40 & [+10.94, +13.85] & +8.22 & [+6.97, +9.66] \\
Unseen drug / 42 & Our harness $-$ direct (1 lab) & +6.38 & [+5.34, +7.53] & +7.55 & [+6.30, +8.92] \\
\midrule
Unseen protein / 42 & Loop $-$ fixed (10 labs) & +0.03 & [-5.53, +4.28] & +4.12 & [-4.10, +8.22] \\
Unseen protein / 42 & Loop $-$ direct (1 lab) & +0.00 & [+0.00, +0.00] & +0.00 & [+0.00, +0.00] \\
Unseen protein / 42 & Loop $-$ direct (10 labs) & +0.00 & [+0.00, +0.00] & +0.00 & [+0.00, +0.00] \\
Unseen protein / 42 & 10 $-$ 1 labs (fixed) & +20.02 & [+0.16, +63.83] & +8.14 & [-8.27, +51.77] \\
Unseen protein / 42 & 10 $-$ 1 labs (direct) & +8.86 & [+2.48, +13.39] & +13.26 & [+3.15, +28.50] \\
Unseen protein / 42 & 10 $-$ 1 labs (loop) & +8.86 & [+2.48, +13.39] & +13.26 & [+3.15, +28.50] \\
Unseen protein / 42 & Our harness $-$ fixed (1 lab) & +20.05 & [+1.84, +61.27] & +12.25 & [-3.01, +53.19] \\
Unseen protein / 42 & Our harness $-$ direct (1 lab) & +8.86 & [+2.48, +13.39] & +13.26 & [+3.15, +28.50] \\
\midrule
Random / 43 & Loop $-$ fixed (10 labs) & +0.54 & [-0.08, +1.11] & +0.44 & [-0.54, +1.44] \\
Random / 43 & Loop $-$ direct (1 lab) & +0.00 & [+0.00, +0.00] & +0.00 & [+0.00, +0.00] \\
Random / 43 & Loop $-$ direct (10 labs) & +0.00 & [+0.00, +0.00] & +0.00 & [+0.00, +0.00] \\
Random / 43 & 10 $-$ 1 labs (fixed) & +7.80 & [+6.82, +8.82] & +7.68 & [+6.32, +9.13] \\
Random / 43 & 10 $-$ 1 labs (direct) & +1.65 & [+0.90, +2.38] & +1.17 & [-0.11, +2.50] \\
Random / 43 & 10 $-$ 1 labs (loop) & +1.65 & [+0.90, +2.38] & +1.17 & [-0.11, +2.50] \\
Random / 43 & Our harness $-$ fixed (1 lab) & +8.35 & [+7.29, +9.37] & +8.12 & [+6.53, +9.76] \\
Random / 43 & Our harness $-$ direct (1 lab) & +1.65 & [+0.90, +2.38] & +1.17 & [-0.11, +2.50] \\
\midrule
Unseen drug / 43 & Loop $-$ fixed (10 labs) & -0.51 & [-1.44, +0.29] & -1.10 & [-2.16, -0.18] \\
Unseen drug / 43 & Loop $-$ direct (1 lab) & +0.00 & [+0.00, +0.00] & +0.00 & [+0.00, +0.00] \\
Unseen drug / 43 & Loop $-$ direct (10 labs) & +0.00 & [+0.00, +0.00] & +0.00 & [+0.00, +0.00] \\
Unseen drug / 43 & 10 $-$ 1 labs (fixed) & +6.68 & [+5.33, +8.25] & +4.66 & [+3.33, +6.36] \\
Unseen drug / 43 & 10 $-$ 1 labs (direct) & -0.15 & [-0.96, +0.67] & -1.43 & [-2.43, -0.42] \\
Unseen drug / 43 & 10 $-$ 1 labs (loop) & -0.15 & [-0.96, +0.67] & -1.43 & [-2.43, -0.42] \\
Unseen drug / 43 & Our harness $-$ fixed (1 lab) & +6.17 & [+4.85, +7.48] & +3.57 & [+2.10, +5.09] \\
Unseen drug / 43 & Our harness $-$ direct (1 lab) & -0.15 & [-0.96, +0.67] & -1.43 & [-2.43, -0.42] \\
\midrule
Unseen protein / 43 & Loop $-$ fixed (10 labs) & -0.96 & [-7.24, +3.07] & -1.83 & [-21.55, +7.54] \\
Unseen protein / 43 & Loop $-$ direct (1 lab) & +0.00 & [+0.00, +0.00] & +0.00 & [+0.00, +0.00] \\
Unseen protein / 43 & Loop $-$ direct (10 labs) & +0.00 & [+0.00, +0.00] & +0.00 & [+0.00, +0.00] \\
Unseen protein / 43 & 10 $-$ 1 labs (fixed) & +4.79 & [-6.61, +28.12] & +0.41 & [-16.48, +31.69] \\
Unseen protein / 43 & 10 $-$ 1 labs (direct) & -1.19 & [-6.22, +6.09] & -6.76 & [-13.95, +6.71] \\
Unseen protein / 43 & 10 $-$ 1 labs (loop) & -1.19 & [-6.22, +6.09] & -6.76 & [-13.95, +6.71] \\
Unseen protein / 43 & Our harness $-$ fixed (1 lab) & +3.83 & [-5.83, +22.32] & -1.42 & [-14.49, +16.88] \\
Unseen protein / 43 & Our harness $-$ direct (1 lab) & -1.19 & [-6.22, +6.09] & -6.76 & [-13.95, +6.71] \\
\midrule
\bottomrule
\end{longtable}
\endgroup

\end{document}
